\documentclass[12pt]{article}
\usepackage[hide-page-numbers]{../../../lecture_notes}
\usepackage{../../../math}
\usepackage{../../../uark_colors}

\hypersetup{
  colorlinks = true,
  allcolors = ozark_mountains,
  breaklinks = true,
  bookmarksopen = true
}

\begin{document}

\begin{center}
  {\Huge\bf Quiz --- Assignment 3}
  
  \smallskip
  {\large\it ECON 5783 --- University of Arkansas}
\end{center}

\vspace{5mm}
\noindent\textbf{Name:} \rule{0.7\textwidth}{0.4pt}

\bigskip
% {\fontsize{20}{24}\selectfont

\begin{enumerate}
  \item Suppose that, conditional on age, smoking during pregnancy is randomly assigned:
  $$
    \left( Y_i(1), Y_i(0) \right) \Perp \text{Mother Smokes}_i \ \mid \ \text{Age}_i.
  $$
  \begin{enumerate}
    \item[(a)] Write the conditional average treatment effect for 30-year-old mothers using observable conditional means.
    Then explain how you would combine age-specific treatment-effect estimates to estimate the overall average treatment effect in the sample.

    \vspace{2\bigskipamount}
    \item[(b)] Give a concrete example of why you think the conditional independence assumption might fail in this setup.
  \end{enumerate}
\end{enumerate}

% }
% 
% \newpage\newpage\ 

\end{document}
